\chapter*{阅读指引}

机器学习理论关心的不是某次训练恰好得到多高的分数，而是有限样本为什么足以支持对未来表现
的判断。回答这个问题，需要同时写清观测与抽样方式、候选决策、损失、算法可见的信息、比较
基准和概率来源。本部分围绕可学习性组织这些对象，并逐步放宽最初的分类与同分布条件。

第一卷已经提供本部分使用的组合维度、集中不等式、随机复杂度、信息量、分布距离和优化
条件；这里不重新主讲这些数学对象，而是说明它们怎样进入学习保证。基础部分提供任务、
模型、损失和范式接口。深度学习理论章只建立推导所需的最小网络形式；网络结构、反向传播
和训练实践由第三卷系统展开。

《机器学习理论基础》建立六章共用的问题设定、PAC标准和归纳偏置，并用二分类例子解释经验
风险与期望风险的区别。随后两章分别研究二分类和多分类中的容量、样本量与计算条件。《一般
学习的可学习性》把零一分类扩展到一般损失，比较一致收敛与算法稳定性；《深度学习的可学习
性》再把这些工具用于过参数化网络，分析参数大小、梯度和参数分布怎样影响选解与泛化。
《迁移学习的可学习性》改变训练与评价的分布关系，说明跨域保证还需要覆盖、共享机制或环境
抽样等额外条件。

六章之间的依赖来自共同的风险、概率量词和可测性约定。基础章规定这些对象，专题章分别改变
标签结构、损失、模型参数化或数据分布；同一个定理只有在这些对象与条件保持一致时才能复用。
可学习性由此说明算法在规定问题中能否用有限经验达到给定标准，不等于模型已经满足真实应用
中的全部可靠性、安全性或责任要求；这些要求由下一部分另行建立。
